\documentclass[../../main.tex]{subfiles}
\begin{document}
\section{The conjecture and why it matters}
\label{sec:kls-orientation}

Part~I is the entry point to this document, and it is meant to be readable on its own. This
section states the Kannan--Lov\'asz--Simonovits conjecture in the forms used below and fixes the
normalizations that make quoted exponents comparable;
Section~\ref{sec:kls-known} records what is known and how the current best bound is built;
Section~\ref{sec:kls-remaining} isolates what is left over, which is the thing every route in
Parts~III and~IV is a response to.

Nothing in Parts~I and~II is a repository result. Both are exposition of the published and
preprint literature, written so that Parts~III--V --- which \emph{are} repository work --- can be
read without leaving the document. Imported statements are attributed individually, and version-1
preprints are marked as such wherever they are used.

The conjecture is a statement about \emph{expansion}. A log-concave measure cannot be cut in two
without paying for the cut, and KLS asserts that in the right normalization the price does not
decay as the dimension grows. Equivalently, in its analytic form, the measure has a spectral gap
bounded below independently of $n$: no function of many variables can have large variance and
small gradient energy at the same time. Both readings are made precise in the next few pages, and
the two-sided comparison \eqref{eq:cheeger-two-sided} is what allows them to be used
interchangeably.

\subsection{The conjecture}
\label{subsec:kls-conjecture}

Let $\mu$ be a log-concave probability measure on $\R^n$. Write $\mu^+(A)$ for the Minkowski
boundary measure of a Borel set $A$, and define the Cheeger constant and its reciprocal, the
inverse Cheeger scale,
\[
  h_\mu=\inf_{A}\frac{\mu^+(A)}{\min\{\mu(A),1-\mu(A)\}},
  \qquad
  \PsiKLS_\mu=h_\mu^{-1}.
\]
The Poincar\'e constant $\CP(\mu)$ is the least constant with
\[
  \Var_\mu(f)\le\CP(\mu)\int_{\R^n}|\nabla f|^2\dd\mu
\]
for every locally Lipschitz $f$.

\begin{warning}[The $\psi$ convention]
\label{rem:psi-convention}
Part of the literature calls $h_\mu$ ``the KLS constant'' and part calls $h_\mu^{-1}$ by the same
name, both writing $\psi_\mu$. This document always writes $h$ for the expansion constant and
$\PsiKLS=h^{-1}$ for the inverse scale, and never uses the bare symbol $\psi$ for either.
Combined with $\CP\asymp\PsiKLS^2$ below, the convention accounts for most apparent
contradictions between quoted exponents: a bound $\PsiKLS_n\lesssim(\log n)^{a}$ is the same
statement as $\CP\lesssim(\log n)^{2a}$, so the same theorem is quoted with two different
exponents depending on which constant the source names.
\end{warning}

For log-concave measures the Cheeger and reverse Cheeger (Buser--Ledoux) inequalities give the
two-sided comparison
\begin{equation}
  \tfrac14\le\frac{\PsiKLS_\mu^2}{\CP(\mu)}\le\pi ,
  \label{eq:cheeger-two-sided}
\end{equation}
so $\CP\asymp\PsiKLS^2$ with universal constants; the explicit constants are recorded in
\cite{Klartag2023Logarithmic,Milman2009Isoperimetric}. The lower bound is Cheeger's inequality in
the normalization $\CP(\mu)\le4h_\mu^{-2}$ used throughout.

\begin{conjecture}[KLS]
\label{conj:kls}\klsstatus{conj:kls}
There is a universal constant $C<\infty$ such that, for every dimension, every isotropic
log-concave $\mu$, and every locally Lipschitz $f$,
\[
  \Var_\mu(f)\le C\int_{\R^n}|\nabla f|^2\dd\mu.
\]
Equivalently up to universal constants, the Cheeger constants of all such measures are bounded
below uniformly in the dimension.
\end{conjecture}

Write
\begin{equation}
  \hstar_n=\inf\bigl\{h_\nu:\ \nu\ \text{isotropic log-concave on }\R^n\bigr\}.
  \label{eq:hstar-def}
\end{equation}
Thus Conjecture~\ref{conj:kls} is equivalently $\inf_n\hstar_n>0$, up to the comparison
\eqref{eq:cheeger-two-sided}.

Three formulations are used interchangeably below.

\begin{enumerate}[label=(\alph*),leftmargin=2.2em]
  \item \textbf{Isotropic form.} For isotropic $\mu$ (that is, $\E_\mu X=0$ and $\Cov_\mu(X)=I_n$),
  \[
    \CP(\mu)=O(1),\qquad \PsiKLS_\mu=O(1),\qquad h_\mu=\Omega(1).
  \]

  \item \textbf{Affine form.} For every log-concave $\mu$,
  \begin{equation}
    \CP(\mu)\le C\norm{\Cov\mu}_\op .
    \label{eq:kls-affine}
  \end{equation}
  The reverse inequality $\CP(\mu)\ge\norm{\Cov\mu}_\op$ is automatic: testing the Poincar\'e
  inequality on the linear function $f(x)=\inner{x}{v}$ with $v$ a top eigenvector of $\Cov\mu$
  gives $\Var_\mu f=\inner{\Cov\mu\,v}{v}$ and $\int|\nabla f|^2\dd\mu=|v|^2$. So
  \eqref{eq:kls-affine} says exactly this: \emph{up to a universal constant, linear functions
  already detect the slowest mode of every log-concave measure.}

  \item \textbf{Geometric form.} The least-expanding arbitrary cut of a convex body is, up to a
  universal factor, no worse than the least-expanding hyperplane cut
  \cite{KannanLovaszSimonovits1995}.
\end{enumerate}

Applying Poincar\'e to $f(x)=|x|^2$ in isotropic position gives
\begin{equation}
  \Var(|X|^2)\le4n\,\CP(\mu),
  \label{eq:kls-implies-thin-shell}
\end{equation}
the implication ``KLS $\Rightarrow$ thin shell'' discussed in
Section~\ref{subsec:kls-solved-neighbours}. There is no known dimension-free converse from this
one radial estimate to KLS. It is therefore accurate to say that thin-shell control alone does
not currently prove KLS, but not that such a converse has been disproved within the log-concave
class.

Finally, no family is known that forces the ratio in \eqref{eq:kls-affine} to grow with $n$. The
entire gap is in the upper bound; it is not a gap between competing upper and lower power laws.

\section{What is known}
\label{sec:kls-known}

Where the problem actually stands: the quantitative frontier, the two neighbouring conjectures
that have been settled and why settling them did not settle KLS, the architecture of the argument
that produces the current best bound, and the exact place the surviving logarithm enters.

\subsection{Status as of 24 August 2026}
\label{subsec:kls-status}

Set
\[
  \PsiKLS_n=\sup_\mu\PsiKLS_\mu,
  \qquad
  C_{\mathrm P,n}=\sup_\mu\CP(\mu),
\]
the suprema over isotropic log-concave measures on $\R^n$. The full conjecture is open. The
status distinction that matters for this repository is between the best \emph{published} bound
and the best \emph{current preprint} claim.

\begin{center}
\begin{tabularx}{\textwidth}{@{}>{\raggedright\arraybackslash}p{0.20\textwidth}>{\raggedright\arraybackslash}p{0.20\textwidth}Y@{}}
\toprule
& Bound & Source and status \\
\midrule
Best published & $\PsiKLS_n\lesssim\sqrt{\log n}$, \quad $C_{\mathrm P,n}\lesssim\log n$
& Klartag's improved Lichnerowicz argument \cite{Klartag2023Logarithmic}. Published; the
benchmark this repository treats as certified literature. \\
\addlinespace
Best current preprint & $\PsiKLS_n\lesssim(\log n)^{1/4}$, \quad $C_{\mathrm P,n}\lesssim\sqrt{\log n}$
& Letwin, arXiv version~1 of 27 July 2026 \cite{Letwin2026QuadraticKLS}. Not peer reviewed; see
the audit in Section~\ref{subsec:mm-audit}. \\
\bottomrule
\end{tabularx}
\end{center}

The quantitative history is the following. Exponents are given for $\PsiKLS_n$ and for
$C_{\mathrm P,n}$ side by side, precisely because of Warning~\ref{rem:psi-convention}.

\begin{center}
\begin{tabularx}{\textwidth}{@{}>{\raggedright\arraybackslash}p{0.13\textwidth}>{\raggedright\arraybackslash}p{0.21\textwidth}>{\raggedright\arraybackslash}p{0.17\textwidth}Y@{}}
\toprule
Date & $\PsiKLS_n$ & $C_{\mathrm P,n}$ & Main development \\
\midrule
1995 & $O(n^{1/2})$ & $O(n)$ & The localization lemma \cite{KannanLovaszSimonovits1995}. \\
\addlinespace
1995--2011 & down to $O(n^{5/12})$ & --- & Bobkov's radial inequality fed by successively better
thin-shell bounds; the last notable pre-localization record \cite{GuedonMilman2011}. \\
\addlinespace
2012/13 & $O(n^{1/3}\sqrt{\log n})$ & $O(n^{2/3}\log n)$ & Stochastic localization
\cite{Eldan2013ThinShell}. \\
\addlinespace
2016; 2024 & $O(n^{1/4})$ & $O(n^{1/2})$ & Fixed Gaussian tilt and the $\Tr(A_t^2)$ stopping
argument \cite{LeeVempala2018,LeeVempala2024}. \\
\addlinespace
2020/21 & $\exp O(\sqrt{\log n\log\log n})$ & same class & Iterated high-Schatten potentials and
affine preconditioning \cite{Chen2021}. \\
\addlinespace
2022 & $O(\log^5 n)$ & $O(\log^{10}n)$ & First polylogarithmic bound; heat flow, spectral
projections, $H^{-1}$ \cite{KlartagLehec2022Polylog}. \\
\addlinespace
2022 & $O(\log^{3.2226}n)$ & $O(\log^{6.4452}n)$ & Two localization representations and
spiked-spectrum estimates \cite{JambulapatiLeeVempala2022KLS}. \\
\addlinespace
2023 & $O(\sqrt{\log n})$ & $O(\log n)$ & Improved Lichnerowicz plus short-time covariance
control \cite{Klartag2023Logarithmic}. \\
\addlinespace
July 2026, v1 & $O((\log n)^{1/4})$ & $O(\sqrt{\log n})$ & Sharp quadratic Poincar\'e inequality,
giving $\kappa_n=O(1)$ \cite{Letwin2026QuadraticKLS}. \\
\bottomrule
\end{tabularx}
\end{center}

\begin{remark}[Reading the table]
\label{rem:history-table-caveats}
Two entries deserve care. The 2016/2024 row is one result: the Lee--Vempala preprint of 2016
appeared in final form in the Annals in 2024, and the repository bibliography records both. The
$O(n^{5/12})$ entry summarizes a sequence of thin-shell improvements rather than a single paper,
and is quoted here only to mark the pre-localization ceiling. Separately, the bibliography
contains \cite{vempala2016kls}, a 2016 announcement of a proof of the full conjecture; it is not
a milestone in this table, and it is listed in the bibliography for completeness of the record
only.
\end{remark}

\subsection{Solved neighbours that are not KLS}
\label{subsec:kls-solved-neighbours}

The classical implication structure is
\begin{equation}
  \mathrm{KLS}\ \Longrightarrow\ \text{thin shell}\ \Longrightarrow\ \text{slicing},
  \label{eq:kls-implication-chain}
\end{equation}
and no reverse implication is known in a dimension-free form. Both of the weaker conjectures have
now moved decisively. Neither resolution proves KLS, and it is worth being precise about why.

\paragraph{Slicing.}
Bourgain's slicing conjecture asks whether the isotropic constants $L_n$ are universally bounded.
Guan proved $L_n\lesssim\log\log n$ and, in the process, obtained the stochastic-localization
trace estimate that turned out to be decisive \cite{Guan2024}. Klartag and Lehec combined that
estimate with $M$-ellipsoids and Shannon--Stam stability to prove $\sup_nL_n<\infty$, published in
2025 \cite{KlartagLehec2025Slicing}; Bizeul subsequently gave an alternative proof through
small-ball estimates \cite{Bizeul2025SmallBallSlicing}. This is a genuine resolution, but slicing
controls determinant and volume information, not the bottom of the full spectrum.

\paragraph{Thin shell.}
For isotropic log-concave $X$ the thin-shell conjecture asks for $\Var(|X|^2)\le Cn$, equivalently
$\E(|X|-\sqrt n)^2\le C$. Klartag and Lehec proved it via parallel couplings of exponential tilts,
nonlinear filtering, optimal transport, $H^{-1}$ estimates, and Guan-type covariance control
\cite{KlartagLehec2025ThinShell}. A July 2026 version-1 preprint of Chen and Klartag sharpened it
to the optimal $\Var(|X|^2)\le8n$, with equality for products of centered exponentials, together
with a sharp third-moment-tensor bound \cite{ChenKlartag2026SharpThinShell}; both statements are
recorded as Theorems~\ref{thm:chen-klartag-thin-shell} and~\ref{thm:chen-klartag-third-moment}.

\paragraph{Why this does not close the gap.}
Eldan's reverse estimate bounds the inverse Cheeger scale by a weighted average of thin-shell
parameters $\sigma_k$ \cite{Eldan2013ThinShell},
\begin{equation}
  \PsiKLS_n\le C\sqrt{(\log n)\sum_{k=1}^n\frac{\sigma_k^2}{k}} .
  \label{eq:eldan-reverse}
\end{equation}
Even with the thin-shell conjecture in hand --- that is, with $\sigma_k=O(1)$ --- the harmonic sum
contributes a $\log n$ and \eqref{eq:eldan-reverse} yields only $\PsiKLS_n\lesssim\log n$, not
$O(1)$. Radial concentration constrains one observable; KLS quantifies over every function and
every measurable cut. The same asymmetry recurs, in sharper form, for Letwin's theorem: it
eliminates every \emph{quadratic} witness, and the first eigenfunction of a general log-concave
diffusion need not be quadratic.

\subsection{The present proof architecture}
\label{subsec:kls-architecture}

Modern progress on KLS is a hybrid argument with four ingredients, no one of which replaces
another. Figure~\ref{fig:kls-architecture} is the shape of the current best bound.

\begin{itemize}[leftmargin=1.6em]
  \item \textbf{Stochastic localization creates Gaussian curvature.} The tilt $e^{-t|x|^2/2}$
  makes the localized measure $t$-strongly log-concave (Section~\ref{sec:family-sl}).
  \item \textbf{Matrix martingale estimates control how much covariance survives.} The covariance
  process $A_t$ obeys an exact Riccati SDE whose noise is the third-moment tensor.
  \item \textbf{Bochner/Lichnerowicz converts curvature plus covariance into a spectral gap.}
  Klartag's improved comparison $\CP\le\sqrt{\norm{\Cov}_\op/t}$
  (Theorem~\ref{thm:improved-lichnerowicz}) is what makes short-time curvature usable.
  \item \textbf{Letwin's preprint supplies the previously missing sharp third-moment estimate},
  through moment maps, Monge--Amp\`ere, Stein kernels, and $H^{-1}$
  (Section~\ref{sec:family-moment-map}).
\end{itemize}

\begin{figure}[htbp]
\centering
\begin{tikzpicture}[
  node distance=8mm and 11mm,
  box/.style={draw,rounded corners=2pt,align=center,font=\small,
              inner sep=4pt,minimum height=8mm,text width=27mm},
  new/.style={box,draw=blue!55!black,thick,fill=blue!6},
  res/.style={box,fill=black!5},
  >={Latex[length=2mm]},
  every path/.style={->,thick,draw=black!60}
]
\node[box] (sl) {Stochastic\\localization};
\node[box,above right=2mm and 14mm of sl] (curv) {Gaussian curvature $t$};
\node[box,below right=2mm and 14mm of sl] (cov) {Covariance process $A_t$};
\node[box,right=of cov] (pot) {Third moments,\\matrix potentials};
\node[new,below=9mm of cov] (mm) {Moment map and\\Monge--Amp\`ere};
\node[new,right=of mm] (qp) {Sharp quadratic\\Poincar\'e};
\node[new,right=of qp] (kap) {$\kappa_n=O(1)$};
\node[box,right=14mm of curv] (il) {Improved\\Lichnerowicz};
\node[res,right=of il] (cp) {$\CP\lesssim\sqrt{\log n}$};
\node[res,right=of cp] (psi) {$\PsiKLS\lesssim(\log n)^{1/4}$};

\draw (sl) -- (curv);
\draw (sl) -- (cov);
\draw (cov) -- (pot);
\draw (mm) -- (qp);
\draw (qp) -- (kap);
\draw (kap) -- (pot);
\draw (curv) -- (il);
\draw (pot) -- (il);
\draw (il) -- (cp);
\draw (cp) -- (psi);
\end{tikzpicture}
\caption{The architecture of the current best bound. Boxed in colour: the July 2026 preprint
contribution. It does not replace stochastic localization; it discharges a crucial input to the
localization--Lichnerowicz pipeline.}
\label{fig:kls-architecture}
\end{figure}

\paragraph{The exact bridge, and where the logarithm lives.}
Define the directional third-moment parameter
\begin{equation}
  \kappa_n=\sup_{\mu,\theta}\bigl\lVert\E_\mu[\inner{X}{\theta}X\otimes X]\bigr\rVert_{\HS},
  \label{eq:kappa-def}
\end{equation}
the supremum over isotropic log-concave $\mu$ on $\R^n$ and $\theta\in S^{n-1}$. The
localization--Lichnerowicz pipeline delivers
\begin{equation}
  \boxed{\ C_{\mathrm P,n}\lesssim\kappa_n\sqrt{\log n}\ }
  \label{eq:kls-bridge}
\end{equation}
(Section~\ref{subsec:sl-where-the-log-lives}), and Letwin's $\kappa_n\le2\sqrt2$
(Proposition~\ref{prop:letwin-kappa}) is what turns \eqref{eq:kls-bridge} into the current record.

The point to retain is \emph{where the surviving $\sqrt{\log n}$ comes from}. It is not hidden in
the moment-map calculation. Klartag--Lehec smooth the maximum covariance eigenvalue by
log-trace-exp,
\[
  \Phi_t=\frac1\beta\log\Tr\bigl(e^{\beta A_t}\bigr),
  \qquad
  \norm{A_t}_\op\le\Phi_t\le\norm{A_t}_\op+\frac{\log n}\beta ,
\]
and approximating $\lmax$ to within a constant forces $\beta\asymp\log n$. The It\^o drift is then
of size $O(\kappa_n^2\log n)$, so covariance control survives only until
$t_*\asymp1/(\kappa_n^2\log n)$, and improved Lichnerowicz converts that time into
$\CP\lesssim t_*^{-1/2}\lesssim\kappa_n\sqrt{\log n}$. The logarithm is the \emph{entropy cost of
replacing a matrix maximum by a soft maximum over $n$ directions}. That diagnosis is what makes
target~3 of Section~\ref{sec:kls-synthesis} concrete.

\section{The shape of the remaining problem}
\label{sec:kls-remaining}

The July 2026 preprints changed the shape of the problem, and it is worth stating the change
precisely rather than as a mood: \emph{the remaining difficulty is no longer quadratic forms or
third moments}. Conditional on \cite{Letwin2026QuadraticKLS}, $\kappa_n=O(1)$ and every quadratic
witness is eliminated. What survives is a difficulty of a different type, and this section states
it in the two forms every route below must answer.

\subsection{Fixed and averaged, against adaptive and uniform}
\label{subsec:kls-adaptive-residue}

Every method surveyed in Part~II controls something. Needles control one-dimensional conditional
measures sharply; stochastic localization controls short-time covariance and, conditionally,
third moments; heat-flow and $H^{-1}$ arguments control coordinate and quadratic spectral mass;
moment maps control fixed deterministic matrix energies $\E\Tr(BHBH)$; parallel coupling controls
linear exponential tilts; Brownian transport controls averaged derivatives to within a
polylogarithm. Section~\ref{subsec:synthesis-table} tabulates these one by one.

What none of them controls is the same object when it is allowed to \emph{adapt}. In every row of
that table the estimate holds for a fixed matrix, a fixed direction, a fixed tilt, or on average
along a path, and KLS needs it uniformly, for an object that may depend on the measure or on the
extremizing function. The gap is not a constant to be improved; it is a quantifier in the wrong
place. A reader who remembers one sentence from Part~I should remember this one, because it is the
sentence each of the four routes of Parts~III and~IV proposes a different way around.

\subsection{The covariance spike, and why the direct repair is false}
\label{subsec:kls-spike-obstruction}

The most natural repair --- bound $\norm{A_t}_\op$ better --- cannot work, and this is not a
guess about difficulty but a proved fact.

\begin{proposition}[Covariance spikes are compatible with KLS]
\label{prop:covariance-spike}\klsstatus{prop:covariance-spike}
Let $\mu$ be the law of $n$ independent centered one-sided exponential coordinates and let $G$ be a
standard Gaussian independent of $X\sim\mu$. Then $\CP(\mu)=O(1)$ by tensorization, yet
\[
  \Prob\!\left(
    \bigl\lVert\Cov(X\mid X+\sqrt sG)\bigr\rVert_\op\ge c_{\rm sp}s
  \right)
  \ge1-\left(1-\tfrac12e^{-s}\right)^n
\]
for a universal $c_{\rm sp}>0$. Consequently, for $s\le\log n$,
\[
  \E\bigl\lVert\Cov(X\mid X+\sqrt sG)\bigr\rVert_\op\gtrsim s
  \qquad\text{and}\qquad
  \Prob\!\left(
    \bigl\lVert\Cov(X\mid X+\sqrt sG)\bigr\rVert_\op\ge c_{\rm sp}s
  \right)\ge1-e^{-1/2}.
\]
\end{proposition}

This is worked out in Section~8.1 and Proposition~65 of Klartag--Lehec \cite{KLnotes}; it is
imported here, not proved.

The consequence is structural, and it is worth stating flatly: \emph{uniform operator-norm control
of the conditional covariance at every scale is false, even for a measure that satisfies KLS.}
Any argument whose only use of the localization path is a pathwise bound on
$\norm{A_t}_\op$ is therefore attempting to prove something stronger than KLS, and something that
is not true. A successful potential must exploit eigenvalue profiles, averaging, tensorization, or
the particular extremizing function --- not better pointwise bounds.

There is a second consequence, and it cuts the other way. Rare covariance spikes can be
\emph{harmless}: the measure above satisfies KLS with a universal constant while spiking. A
potential that charges the full largest eigenvalue whenever a spike occurs is therefore
overcharging, and overcharging in a way that is easy to detect.

\subsection{The tensorization test}
\label{subsec:kls-tensorization-test}

That gives a cheap and surprisingly discriminating test, which any proposal --- including every
route below --- should be made to pass before it is developed.

\begin{quote}
  \emph{Evaluate the proposed potential on a product of $n$ independent copies of a
  one-dimensional measure. A quantity that is not tensorization-aware will charge $n$ independent
  coordinates $n$ times for a phenomenon that costs $O(1)$.}
\end{quote}

The test is not a formality. Section~\ref{sec:product-stress} applies it to the all-cut Carleson
estimate of Route~E and it is where the rank-one candidate is refuted; the same computation is
what identifies the surviving adapted alignment problem. A proposal that fails it is not merely
inefficient, it is false at scale, and the failure is usually visible in a page.

Every live route of this repository is a response to the two constraints of this section:

\begin{itemize}[leftmargin=1.6em]
  \item Route~E (Section~\ref{sec:introduction}) retains the two-color covariance of one fixed
  candidate cut, rather than the full spectrum of $A_t$.
  \item Route~S (Section~\ref{sec:spectral-route}) retains the covariance tensor of one fixed
  first eigenfunction.
  \item Route~C (Section~\ref{sec:moment-map-cmh}) abandons stochastic localization altogether
  and works in deterministic moment-map coordinates.
  \item Route~F (Section~\ref{sec:conditional-fiber-frame}) replaces Euclidean directions by one
  test-independent isotropic frame of inverse-variance-normalized conditional line resamplings.
\end{itemize}

Section~\ref{sec:frontier-atlas} compares them side by side, with the advance and the exact
blocker of each; the chapters of Parts~III and~IV then develop them at their natural lengths.

\end{document}
